\subsection{A Heterogeneity Score to Complement Performance Portability Metrics}
\label{sec:heterogeneity_score}

As \citet{yokota_beginners_2018} already stated, it would be desirable to not only assess the performance portability of some platform set but to get a feeling for whether said platform set is representable for the platforms one is interested in. 
For example, if we have an application supporting \glspl{gpu} and \glspl{cpu} where the performance efficiency of \glspl{gpu} is noticeably higher than for \glspl{cpu}, we can improve the performance portability score by adding more \glspl{gpu} or removing the \glspl{cpu}. 
However, having significantly more \glspl{gpu} in our platform set than \glspl{cpu} may be bad if our ultimate goal is to also support \glspl{cpu} as well as \glspl{gpu}. 
Therefore, given the architectural classes our application should support, it would be advisable to have some heterogeneity score that grades the representability of our platform set. 
Such a score can help us to identify whether our selected set of platforms is representative and whether the conclusions we want to draw with \pp or \mpp are correct. 
For example, if our performance portability score is close to $1$ but our heterogeneity score is close to $0$, that would indicate that our application may be performance portable according to \pp and \mpp, but our analysis misses many hardware platforms that are of interest to us, possibly skewing the results. 
In other words, can we conclude from \pp and \mpp alone that our application is performance portable across the hardware platforms of interest, or is it only performance portable due to cherry-picking the best hardware platforms? 
On the other hand, if \pp or \mpp respectively are close to $1$ and our heterogeneity score is also close to $1$, we could conclude that our selected set of platforms represents our platforms of interest in such a good way that we can indeed say that our application is performance portable \textit{with respect to the architectural classes we want to support}.

Based on this observation, I propose the \textit{Heterogeneity Score} $HS$ based on the Shannon Entropy~\cite{shannon_mathematical_1948} that can go hand in hand with the performance portability metrics to show whether we can really draw conclusions from our \pp or \mpp scores. 

\begin{definition}[Heterogeneity Score]\label{def:hetero_score}
Given different architectural classes $\mathbf{C}$, and the set of all platforms $\mathbf{H}$ where each $\mathbf{p \in H}$ is associated with an architectural class $\mathbf{c \in C}$, let the heterogeneity score be defined as:

$$HS(C, H) = \begin{cases}
                \frac{-\sum\limits_{c \in C; p_c > 0} p_c \log_2 p_c}{\log_2\abs{C}} & \text{if } \abs{C} > 1 \\
                0                                                                    & \text{otherwise}
              \end{cases}$$

where $p_c = \frac{\abs{\{ p \in H \mid p \text{ is associated with } c \}}}{\abs{H}}$ is the ratio of the number of platforms in the current architectural class $c$ relative to the total number of platforms. 
\end{definition}
Note that we only sum over all $p_c$ that are greater than zero; otherwise, the logarithm would be undefined. 

If we only have a single architectural class, our heterogeneity score is defined to be zero. 
This aligns with the intended usage of the $HS$ score and a performance portability metric: if we only use a single architectural class, are we really talking about portability? 
In my opinion, no. 
Likewise, if the number of platforms for each architectural class is the same, $HS$ becomes one since we have perfect heterogeneity across all hardware platforms we are interested in. 
The further away from an optimal distribution or the fewer architectural classes of interest we use, the smaller our heterogeneity scores get. 
For example, if we have two architectural classes, $c_1$ and $c_2$, and currently three hardware platforms, one associated with the class $c_1$ and two with the class $c_2$, we would have a heterogeneity score of $\num{0.918}$. 
If we add a new hardware platform associated with the architectural class $c_2$, our heterogeneity score reduces to $\num{0.811}$. 
This aligns with the observations of \citet{yokota_beginners_2018}, which showed that we can improve the scores of \pp or \mpp by adding more hardware platforms that are already well optimized. 
If we combine these \pp and \mpp scores now with the new $HS$ metric, we can directly see that we did not improve the performance portability in a meaningful way if we are interested in portability. 

Note that it is up to the user to define the granularity of the different architectural classes. 
For example, we can define \glspl{gpu} and \glspl{cpu} as two classes. 
Likewise, we can also define more fine-grained classes by splitting the \glspl{gpu} according to their vendors, creating unique architectural classes for \glspl{gpu} from NVIDIA, AMD, or Intel. 
The same holds for \glspl{cpu}, where we can define separate classes for different \glspl{cpu} micro-architectures, like x86-64, ARM, or POWER. 
This flexibility allows us to adjust the heterogeneity score to our needs. 

Therefore, I think the proposed heterogeneity score $HS$ is a great addition to the well-established performance portability metrics, which may help to better understand the performance portability of one's application. 
